\section{Analysis 1: Liquidity Threshold Effect}
\label{sec:analysis1}

[SECTION TO BE EXPANDED BY OPUS]

% When filled in:
% - Hypothesis 1 restatement
% - Methodology (stratification by volume bins)
% - Main finding: ~200-trade inflection point
% - Multiple TikZ charts showing threshold effect
% - Interpretation and robustness

\subsection{Methodology}
\label{subsec:a1_method}

We stratify the 7.3M market sample by trading volume into five bins...

\subsection{Main Finding: The $\sim$200-Trade Threshold}
\label{subsec:a1_finding}

% TikZ figure showing threshold effect
\begin{figure}[h]
\centering
\begin{tikzpicture}
  \begin{axis}[
    title={MAE vs. Trading Volume (Log Scale)},
    xlabel={Trading Volume (log scale)},
    ylabel={Mean Absolute Error (\%)},
    width=0.9\textwidth,
    height=0.6\textwidth,
    xmode=log,
    grid=major,
    legend pos=upper right
  ]
    % Placeholder data - will be filled in by Opus
    \addplot[color=darkblue, mark=o, thick] coordinates {
      (10, 18)
      (50, 16.5)
      (100, 16.37)
      (250, 12.0)
      (500, 9.13)
      (2000, 7.29)
      (10000, 3.87)
    };
    \addlegend{Observed MAE}
    
    % Add threshold indicator
    \draw[dashed, red, thick] (200, 0) -- (200, 25);
    \node[anchor=north, red] at (200, 2) {$\sim$200 trades};
  \end{axis}
\end{tikzpicture}
\caption{Sharp inflection point in calibration error at approximately 200 trades. Below this threshold, each additional trade reduces MAE by $\sim$0.05\%; above it, marginal improvement flattens.}
\label{fig:threshold_main}
\end{figure}

\subsection{Robustness Checks}
\label{subsec:a1_robust}

% Additional charts and discussion
